\section{Task State, Trees, and Local Laws}\label{sec:min-framework}
\label{sec:min-manifesto-state}

The algebra is small. At each node of the tree we want a summary the
oracle cannot distinguish from the raw span beneath it, and three local
laws say how that property propagates from leaves to root. Whatever the
summarizer is, the local laws, the audit, and the gap bound treat it the
same way.

\subsection{Two Pieces a Researcher Chooses: $f$ and $g$}\label{sec:min-fg-pieces}

Let $\Strings$ be finite token sequences with concatenation
$u \concat v$. The two pieces a researcher chooses are the summarizer
$g:\Strings\rightsquigarrow\Strings$ and the readout
$f:\Strings\rightsquigarrow\Y$. The summarizer compresses each leaf
or merge call into a stored summary; the readout scores the root summary.
The pair is an approximate sufficient statistic for the task on a
tree merge: knowing $g(x)$ should determine $f(x)$ up to a controlled
error, and the same identity should hold after pairwise merges.

The pair can be hand-crafted in closed form, as in the classical
mergeable-sketch line, or learned in a parameterized function class.
C-TreePO is the same algebra under the second choice, with the
particular function class fixed by the application: prompt or weight
optimization against the oracle for LLM-shaped $g$, a neural operator
\citep{Kovachki23NeuralOperator} for continuous structured inputs, or
a fine-tuned small model for distillation cascades. The Manifesto
application below uses LLM summaries and LLM scores: $g$ is a
prompted Gemma summarizer applied at every leaf and merge, and $f$ is
a prompted Gemma scorer that returns a seven-point policy score. The
framework runs unchanged across instantiations along the lineage:

\begin{itemize}[leftmargin=1.5em]
  \item \textbf{Classical mergeable sketches.} For Count-Min
    \citep{CormodeMuthukrishnan2005}, $g$ hashes each item into one
    counter per row and merges by element-wise counter sum; $f$
    returns the row-minimum for a queried item. The pair is an exact
    sufficient statistic for item frequency on a fixed query catalog.
    HyperLogLog (cardinality) and Bloom filters (membership) fit the
    same template with different hand-designed states
    (Apps.~\ref{app:min-hll}, \ref{app:min-algebra}).
  \item \textbf{Learned neural operators.} $g$ is a Fourier neural
    operator and $f$ is a small MLP readout. The Markov changepoint task
    is the worked example (Apps.~\ref{app:min-neural},
    \ref{app:min-markov}).
  \item \textbf{Prompted LLMs (the case studied here).} $g$ is a
    prompted LLM summarizer, $f$ is a prompted LLM scorer; the
    Manifesto results in Section~\ref{sec:min-manifesto-results}.
  \item \textbf{Distillation cascade (forward-looking).} $g$ is a small
    fine-tuned LM trained on node-level labels from a stronger proxy;
    $f$ is a light scorer trained alongside it
    (Section~\ref{sec:min-discussion}).
\end{itemize}

The local laws, the audit, and the gap bound do not see which $(f, g)$
pair was chosen.

\subsection{Running Example: Economic Policy Position}\label{sec:min-running-example}

The running example throughout is a manifesto's economic-policy
position on the public-services-versus-taxation scale used in the
Manifesto Project benchmark \citep{BenoitEtAl2025,ManifestoProject2025a}.
The trusted oracle $\fstar$ is the seven-point score a rubric-trained
coder would assign to the full text. The summarizer $g$ is asked to
preserve the rubric-relevant evidence in any span: commitments to
public services, taxation, welfare, market intervention, privatization,
fiscal restraint, and the qualifications that determine how those
claims should be read. The readout $f$ scores a (raw or summarized)
string on the same seven-point scale, so $f\circ g$ is what the
deployed pipeline computes on a leaf or merged summary. Two texts can
look different and sit in the same fiber of $\fstar$ if they carry the
same economic evidence. Two texts can look similar and sit in
different fibers if one contains a fiscal qualification the other
omits.

\subsection{The Tree}\label{sec:min-tree}

A document $x$ is partitioned into contiguous leaves
$(b_1,\ldots,b_k)$, and those leaves are attached to a binary tree.
Both pieces of structure are application choices. The partition can
be uniform chunks of fixed size, paragraph boundaries, section breaks,
or any other contiguous segmentation. The binary merge tree built on
top of the partition is similarly free; under the local laws below,
any two merge trees on the same leaves give the same expected root
value (Corollary~\ref{cor:schedule}, schedule invariance).
Section~\ref{sec:min-manifesto-results} reports the empirical
extension: across a wide range of partition choices on the Manifesto
Project benchmark, the root estimand stays stable as well.

Each leaf stores $s_i = g(b_i)$. Each internal node $u$ stores
$s_u = g(s_{u_L}\concat s_{u_R})$. Only stored summaries move through
the pipeline at deployment; the raw span beneath a node, written
conceptually as $S(u)$, is retained for audit.

\begin{figure}[t]
    \centering
    \includegraphics[width=0.92\linewidth]{01_base_plain.pdf}
    \caption{\textbf{A compression tree reduces a document to one root
    summary by recursively merging adjacent child summaries.} In the
    Manifesto application, every stored summary stands in for the
    policy-evidence state of the span beneath it.}
    \label{fig:min-plain-tree}
\end{figure}

\subsection{The Three Local Laws}\label{sec:min-laws}

\begin{definition}[Oracle equivalence]\label{def:oracle-equivalence}
For strings $u,v\in\Strings$ and oracle $\fstar:\Strings\rightsquigarrow\Y$
with metric $\metric$ on $\Y$,
\[
  u\sim_{\fstar}v
  \iff
  \metric\big(\fstar(u),\fstar(v)\big)=0.
\]
We write $u\sim v$ for $u\sim_{\fstar}v$ when the oracle is clear from
context. The relation is the equivalence kernel of $\fstar$, and its
fibers are the level sets $\{x:\fstar(x)=\fstar(x_0)\}$. Expectations
of distances under $\fstar$ are taken with respect to $\sim_{\fstar}$.
\end{definition}

Preservation is always relative to this task relation. A valid summary
need not preserve style, rhetoric, or facts irrelevant to the
seven-point score. It must preserve whatever can move $\fstar$.

\begin{assumption}[Context-compatible oracle]\label{ass:context}
For every context $w\in\Strings$, $u\sim_{\fstar}v$ implies both
$w\concat u\sim_{\fstar}w\concat v$ and
$u\concat w\sim_{\fstar}v\concat w$.
\end{assumption}

This is an application assumption. It is plausible for content rubrics
whose evidence can be localized and substituted. It can fail for
targets like tone, irony, or persuasion, where a locally equivalent
span may change how neighboring text is read.

\paragraph{C1: Sufficiency.}
\begin{equation}\tag{C1}\label{eq:leaf_sufficiency}
  \forall\text{ realized leaf }b:\qquad g(b)\sim_{\fstar}b.
\end{equation}
A leaf summary keeps the evidence the oracle would read in its source
span.

\paragraph{C2: Idempotence.}
\begin{equation}\tag{C2}\label{eq:leaf_idempotence}
  \forall\text{ stored summary }s\in\operatorname{range}(g):\qquad
  g(s)\sim_{\fstar}s.
\end{equation}
Re-summarizing an already-reduced string does not shift the task value.

\paragraph{C3: Merge consistency.}
\begin{equation}\tag{C3}\label{eq:leaf_compatibility}
  \forall u,v\in\Strings:\qquad
  u\concat v
  \sim_{\fstar}g(u\concat v)
  \sim_{\fstar}g\!\left(g(u)\concat g(v)\right).
\end{equation}
Summarizing children separately and then merging preserves the same
cross-span interaction as summarizing the raw union at once.

\subsection{From Hand-Crafted to Learned Mergeable Sketches}\label{sec:min-sketch-reduction}

Classical mergeable sketches and C-TreePO instantiate the same
encode-merge-query template. What differs is who designs $(f, g)$ and
what guarantees they carry. In the classical case $(f, g)$ is
hand-crafted to make an exact sufficiency proof go through on a fixed
query catalog: Count-Min picks a counter-array state, an element-wise
sum merge, and a row-min readout for item frequency; HyperLogLog
picks register-wise max and a cardinality estimator that handles the
harder case where $|A|$ and $|B|$ alone do not determine
$|A\cup B|$. The state can carry strictly more information than the
per-query answer, which is why sketches compose \emph{state} and not
scalars.

C-TreePO keeps the encode-merge-query shape and replaces hand-crafting
with optimization. $(f, g)$ live in a parameterized function class
(a prompted or weight-tuned LLM, a neural operator, a fine-tuned small
model), and the closed-form sufficiency proof becomes an audit over
residual local-law violations on the realized calls of the deployed
tree. For the manifesto application, the state is a learned or
prompted surrogate for rubric-relevant evidence; for continuous
structured-input tasks (the Markov walkthrough in
App.~\ref{app:min-markov}), the state is the latent function the
neural operator learns. The local-law machinery does not see which
function class was chosen.

\begin{theorem}[Local laws $\Leftrightarrow$ mergeable-summary algebra]\label{thm:sketch-equivalence}
The following are equivalent for an oracle $\fstar$ and a tree
operator $(g, f)$:
\begin{enumerate}[leftmargin=1.5em]
  \item $(g, f)$ satisfies C1, C2, and C3 of
    Definition~\ref{def:oracle-equivalence};
  \item there is a mergeable-summary scheme
    $(\mathrm{encode}, \oplus, \mathrm{query})$ on a state space
    $\mathcal{S}$ with a decoder $\mathrm{decode}:\mathcal{S}\to\Strings$
    such that
    $g(x)=\mathrm{decode}(\mathrm{encode}(x))$,
    $\mathrm{decode}(\mathrm{encode}(s_L)\oplus\mathrm{encode}(s_R))=g(s_L\concat s_R)$, and
    $f=\mathrm{query}\circ\mathrm{encode}$.
\end{enumerate}
The pointwise approximate version transports one-to-one: nodewise
$\varepsilon$-budgets at C1/C3 correspond to leaf-preservation and
merge-compatibility budgets on the sketch side, and conversely.
\end{theorem}

The reverse construction takes $\mathcal{S}=\operatorname{range}(g)$
with $\mathrm{encode}=g$, $\oplus(s,t)=g(s\concat t)$, and
$\mathrm{query}=f$; C3 supplies the mergeable axiom directly. The
proof of both directions, including the approximate transport, is in
Appendix~\ref{app:min-algebra}; the Lean development verifies the
forward direction (\texttt{local\_laws\_of\_sketch} in
\texttt{SketchRecovery.lean}) and its approximate variant
(\texttt{approx\_bundle\_of\_sketch}). The C-TreePO theorem stack
(root preservation, schedule invariance, DPO/GRPO equivalence) lifts
to any mergeable sketch through this equivalence, and a learned
summarizer that satisfies the audit inherits the same closed-form
guarantees.

\subsection{What LLM-Shaped $g$ Adds: Observability}\label{sec:min-observability}

Classical sketch state is a register tuple. It is inspectable
algebraically (the sufficiency proof goes through on paper) and opaque
to a reader: an HLL register array does not tell anyone which items
were hashed into it. LLM-shaped $g(x)$ is text. It is inspectable
directly, by a reader or by another LLM, in the same units the leaf
$x$ already lives in.

Two consequences carry through to the audit machinery
(Section~\ref{sec:min-audit}). First, the C2 idempotence check has a
literal reading: a sampled stored summary $s$ either preserves the
load-bearing economic evidence under the rubric or it does not, and a
gold-standard coder can score it on the same rubric used at the
document level. Second, sufficiency on a fiber can be probed by
counterfactual sampling. Pairs $(x, x')$ with $g(x)\approx g(x')$
under a textual similarity check, audited for
$\fstar(x)\approx\fstar(x')$, give a sufficiency diagnostic with no
gradient information. Classical sketches have no analogue; their
fibers are not in human-comparable units.

The summary itself becomes a supervision target alongside the
downstream loss. Distillation against teacher summaries, stylistic
constraints on the summary text, and rubric-shaped scaffolding all act
on $g$ directly rather than on $f\circ g$
(Section~\ref{sec:min-distillation}).

The laws are stated globally; the framework only ever needs them on the
realized calls of a deployed tree. Whether they hold often enough on
those calls is not something one asserts. It is something the audit
estimates, and Section~\ref{sec:min-audit} is how.
